\introsection{Объект и предмет исследования}

Объектом исследования являются процессы взаимодействия пользователей с реляционными базами данных в задачах аналитики больших данных. Предметом исследования выступают методы и алгоритмы преобразования естественного языка в SQL-запросы, а также механизмы повышения точности и надёжности их генерации с использованием больших языковых моделей.

\introsection{Теоретические и методические основы исследования}

Теоретическую основу работы составляют методы обработки естественного языка, архитектуры трансформеров, подходы семантического поиска и принципы построения систем класса Retrieval-Augmented Generation. В качестве методической базы использованы алгоритмы векторизации текста, методы измерения семантической близости (косинусное сходство), а также агентные подходы (ReAct) для реализации многошагового взаимодействия с внешними инструментами. Для повышения качества генерации применён метод итеративной самокоррекции, основанный на анализе ошибок выполнения SQL-запросов.

\introsection{Новизна и основные результаты исследования}

В рамках выполнения работы были получены результаты, непосредственно связанные с поставленными задачами. Проведён анализ современных методов Text-to-SQL, на основе которого обоснован выбор архитектуры решения. Спроектирована и реализована RAG-архитектура, обеспечивающая эффективный поиск контекста по структуре базы данных. Разработан механизм самокоррекции SQL-запросов, позволяющий повысить точность генерации за счёт использования обратной связи от СУБД. Проведена экспериментальная оценка качества моделей, показавшая рост метрик выполнения при использовании предложенных методов. Реализована микросервисная архитектура системы, обеспечивающая масштабируемость и модульность решения.

\introsection{Апробация и внедрение результатов}

Результаты исследования были апробированы в рамках экспериментального тестирования на демонстрационной базе данных «Авиаперевозки». Разработанный прототип системы использовался для выполнения набора тестовых SQL-задач, что позволило оценить эффективность предложенных алгоритмов и подтвердить их практическую применимость.

\introsection{Библиографическое описание}

Выпускная квалификационная работа посвящена разработке интеллектуального модуля для интерактивного анализа данных в реляционных базах данных на основе больших языковых моделей. В работе рассмотрены современные подходы к задаче Text-to-SQL, предложены методы повышения точности генерации запросов и реализована программная система, обеспечивающая взаимодействие пользователя с базой данных на естественном языке.
